% Generated draft: requires mathematical and source review.
\documentclass[11pt]{article}
\usepackage[margin=25mm]{geometry}
\usepackage{fontspec}
\setmainfont{Latin Modern Roman}
\usepackage{amsmath,amssymb,mathtools,unicode-math}
\setmathfont{Latin Modern Math}
\usepackage{enumitem,needspace,xurl,hyperref}
\setlength{\parindent}{0pt}
\setlength{\parskip}{5pt}
\setlength{\emergencystretch}{3em}
\newcommand{\N}{\mathbb{N}}
\newcommand{\Z}{\mathbb{Z}}
\newcommand{\Q}{\mathbb{Q}}
\newcommand{\R}{\mathbb{R}}
\newcommand{\C}{\mathbb{C}}
\newcommand{\F}{\mathbb{F}}
\newcommand{\A}{\mathbb{A}}
\newcommand{\PP}{\mathbb{P}}
\newcommand{\E}{\mathbb{E}}
\newcommand{\eps}{\varepsilon}
\newcommand{\dd}{\,\mathrm{d}}
\begin{document}
\section*{UnsolvedMath ID: 30002554}
\textbf{Source problem code:} catalogue entry
\subsection{Definitions and mathematical statement}
Let \(\mathcal H(4)\) denote the stratum of pairs \((X,\omega)\), where \(X\) is a compact connected Riemann surface of genus \(3\) and \(\omega\) is a nonzero holomorphic one-form with exactly one zero, of order \(4\). A square-tiled surface is such a translation surface admitting a finite translation covering
\[
 p:X\longrightarrow \C/(\Z+i\Z),\qquad \omega=p^*(\dd z),
\]
ramified only over the origin. No primitivity assumption on this covering is supplied.

The Veech group of \((X,\omega)\) is the group of derivatives of its orientation-preserving affine automorphisms in translation coordinates. To discuss congruence subgroups, an integral realization must be fixed. For example, choosing a positively oriented basis of the absolute period lattice
\[
\Lambda_\omega=\left\{\int_\gamma\omega:\gamma\in H_1(X,\Z)\right\}\subset\C
\]
realizes the Veech group as a finite-index subgroup \(\Gamma(X,\omega)\leq\operatorname{SL}_2(\Z)\). This choice is relevant because the source does not specify a normalization for nonprimitive square-tiled surfaces. For an integer \(N\geq1\), define the principal congruence subgroup
\[
\Gamma(N)=\ker\!\left(\operatorname{SL}_2(\Z)\longrightarrow\operatorname{SL}_2(\Z/N\Z)\right).
\]
A subgroup \(\Gamma\leq\operatorname{SL}_2(\Z)\) is a congruence subgroup if \(\Gamma(N)\subseteq\Gamma\) for some \(N\geq1\).

Under the standard interpretation, the Prym locus in \(\mathcal H(4)\) consists of pairs admitting a holomorphic involution \(\tau:X\to X\) such that
\[
\tau^*\omega=-\omega,
\qquad \operatorname{genus}(X/\langle\tau\rangle)=1.
\]
The supplied record does not explicitly define its intended Prym locus; this interpretation therefore requires confirmation against the source.

\textbf{Source question.} Determine how far the Veech groups of square-tiled surfaces in \(\mathcal H(4)\), with particular attention to those in the Prym locus, are from being congruence subgroups. The question concerns the full collection of such surfaces, without a supplied bound on the number of squares. The phrase ``how far'' does not specify a numerical invariant, comparison criterion, or desired classification. Consequently, the record does not determine a unique precise quantitative statement; the definition of congruence above specifies the reference property but does not resolve the intended notion of congruence deficiency.

The supplied dated assessment labels the problem partially solved and attributes precise deficiency results for origamis in \(\mathcal H(2)\) and for the Prym locus in \(\mathcal H(4)\) to Weitze--Schmithusen. It supplies neither the statements nor hypotheses of those results and supplies no classification for all origamis. These are reported curation claims, not independently verified mathematical conclusions.
\subsection{Short English statement}
The question asks how the Veech groups of square-tiled genus-three translation surfaces with one zero of order four fail to be congruence groups, especially on the Prym locus. The intended measure of failure, normalization, and exact Prym-locus convention remain unspecified. The supplied partially solved label and cited partial results require independent review.
\subsection{Sources}
\begin{itemize}
\item Local UnsolvedMath registry, \texttt{data/unsolvedmath/problems.json}, permanent ID \(30002554\), source problem code \texttt{catalogue entry}, ``Congruence Deficiency of Veech Groups in \(\mathcal H(4)\).'' The supplied record gives proposed year 2014 and names no individual proposer.
\item \emph{Flat Surfaces and Dynamics on Moduli Space} (2014), Oberwolfach Reports, pp.~869--941, DOI 10.4171/owr/2014/15. Supplied workshop-report provenance and literature-source link: \url{https://doi.org/10.4171/owr/2014/15}.
\item Weitze--Schmithusen, IMRN, DOI 10.1093/imrn/rnt268. Supplied literature-source link: \url{https://doi.org/10.1093/imrn/rnt268}. The record attributes deficiency results in \(\mathcal H(2)\) and the Prym locus in \(\mathcal H(4)\) to this reference, without supplying their exact statements.
\item Literature assessment date: \texttt{2026-08-21}. The supplied dated curation labels the question partially solved and states that no classification for all origamis is supplied. This assessment derives from the supplied Oberwolfach curation and has not been independently verified here; its status and cited claims require independent checking.
\item The registry labels the statement \texttt{corrected\_verified} and reports a manual semantic, typographic, and LaTeX review. This is supplied editorial metadata, not independent verification of the formulation or literature claims.
\end{itemize}
\end{document}
